\chapter{Thiết kế Cấu trúc Dữ liệu và Giải thuật Cây BST}
\label{chap:chuong2}

\textit{Phụ trách biên soạn chính: Phan Thế Hiệp (M2 - BST và Thuật toán).}

\section{Tổng quan về Cây Nhị phân Tìm kiếm (BST)}
\subsection{Định nghĩa toán học và Tính chất}
Cây nhị phân tìm kiếm $T$ là một cây nhị phân hữu hạn, trong đó mỗi node chứa một khóa $k$ thỏa mãn điều kiện:
\begin{itemize}
    \item Mọi khóa trong cây con trái của node $u$ đều nhỏ hơn khóa tại $u$: $\forall x \in LeftSubtree(u), key(x) < key(u)$.
    \item Mọi khóa trong cây con phải của node $u$ đều lớn hơn khóa tại $u$: $\forall y \in RightSubtree(u), key(y) > key(u)$.
    \item Cả cây con trái và cây con phải cũng là các cây nhị phân tìm kiếm.
\end{itemize}

\section{Thiết kế Cấu trúc Dữ liệu Sách (Book) và Node Cây (BSTNode)}
\subsection{Cấu trúc thực thể Book}
Khóa chính so sánh là chuỗi \texttt{ISBN}.
\subsection{Cấu trúc Node cây (BSTNode)}
Node lưu trữ đối tượng \texttt{Book} và hai con trỏ \texttt{left}, \texttt{right}.

\section{Các Giải thuật Thao tác Cốt lõi}
\subsection{Giải thuật Chèn (Insert)}
\subsection{Giải thuật Tìm kiếm nhị phân (Search)}
\subsection{Giải thuật Xóa Node (Erase) với 3 trường hợp}
\subsection{Các giải thuật Duyệt cây (NLR, LNR, LRN)}
\subsection{Giải thuật Giải phóng Bộ nhớ đệ quy (Clear / Destructor)}
